\documentclass{ximera}

\begin{document}
	\author{Stitz-Zeager}
	\xmtitle{Exercises for Par Frac}{}

\mfpicnumber{1} \opengraphsfile{ExercisesforParFrac} % mfpic settings added 


\label{ExercisesforParFrac}

\begin{question}
In Exercises \ref{parfracformfirst} - \ref{parfracformlast},  find only the \emph{form} needed to begin the process of partial fraction decomposition.  Do not create the system of linear equations or attempt to find the actual decomposition.

\begin{problem}\label{parfracformfirst}
$\frac{7}{(x - 3)(x + 5)}$

\begin{solution}
$\frac{A}{x - 3} + \frac{B}{x + 5}$
\end{solution}
\end{problem}

\begin{problem}
$\frac{5x + 4}{x(x - 2)(2 - x)}$

\begin{solution}
$\frac{A}{x} + \frac{B}{x - 2} + \frac{C}{(x - 2)^{2}}$
\end{solution}
\end{problem}

\begin{problem}
$\frac{m}{(7x - 6)(x^{2} + 9)}$

\begin{solution}
$\frac{A}{7x - 6} + \frac{Bx + C}{x^{2} + 9}$
\end{solution}
\end{problem}

\begin{problem}
$\frac{ax^{2} + bx + c}{x^3(5x + 9)(3x^{2} + 7x + 9)}$

\begin{solution}
$\frac{A}{x} + \frac{B}{x^{2}} + \frac{C}{x^{3}} + \frac{D}{5x + 9} + \frac{Ex + F}{3x^{2} + 7x + 9}$
\end{solution}
\end{problem}

\begin{problem}
$\frac{\text{A polynomial of degree } < 9}{(x + 4)^{5}(x^{2} + 1)^{2}}$

\begin{solution}
$\frac{A}{x + 4} + \frac{B}{(x + 4)^{2}} + \frac{C}{(x + 4)^{3}} + \frac{D}{(x + 4)^{4}} + \frac{E}{(x + 4)^{5}} + \frac{Fx + G}{x^{2} + 1} + \frac{Hx + I}{(x^{2} + 1)^{2}}$
\end{solution}
\end{problem}

\begin{problem}\label{parfracformlast}
$\frac{\text{A polynomial of degree } < 7}{x(4x - 1)^{2}(x^{2} + 5)(9x^{2} + 16)}$

\begin{solution}
$\frac{A}{x} + \frac{B}{4x - 1} + \frac{C}{(4x - 1)^{2}} + \frac{Dx + E}{x^{2} + 5} + \frac{Fx + G}{9x^{2} + 16}$
\end{solution}
\end{problem}

\end{question}

\begin{question}
In Exercises \ref{findparfracfirst} - \ref{findparfraclast},  find the partial fraction decomposition of the following rational expressions.

\begin{problem}\label{findparfracfirst}
$\frac{2x}{x^{2} - 1} = \frac{\answer{1}}{x + 1} + \frac{\answer{1}}{x - 1}$
\end{problem}

\begin{problem}
$\frac{-7x + 43}{3x^{2} + 19x - 14}$

\begin{solution}
$\frac{-7x + 43}{3x^{2} + 19x - 14}= \frac{5}{3x - 2} - \frac{4}{x + 7}$
\end{solution}
\end{problem}

\begin{problem}
$\frac{11z^{2} - 5z - 10}{5z^{3} - 5z^{2}}$

\begin{solution}
$\frac{11z^{2} - 5z - 10}{5z^{3} - 5z^{2}} = \frac{3}{z} + \frac{2}{z^{2}} - \frac{4}{5(z - 1)}$
\end{solution}
\end{problem}

\begin{problem}
$\frac{-2z^{2} + 20z - 68}{z^{3} + 4z^{2} + 4z + 16}$

\begin{solution}
$\frac{-2z^{2} + 20z - 68}{z^{3} + 4z^{2} + 4z + 16} = -\frac{9}{z + 4} + \frac{7z - 8}{z^{2} + 4}$
\end{solution}
\end{problem}

\begin{problem}
$\frac{-s^{2} + 15}{4s^{4} + 40s^{2} + 36}$

\begin{solution}
$\frac{-s^{2} + 15}{4s^{4} + 40s^{2} + 36} = \frac{1}{2(s^{2} + 1)} - \frac{3}{4(s^{2} + 9)}$
\end{solution}
\end{problem}

\begin{problem}
$\frac{-21s^{2} + s - 16}{3s^{3} + 4s^{2} - 3s + 2}$

\begin{solution}
$\frac{-21s^{2} + s - 16}{3s^{3} + 4s^{2} - 3s + 2} = -\frac{6}{s + 2} - \frac{3s + 5}{3s^{2} - 2s + 1}$
\end{solution}
\end{problem}

\begin{problem}
$\frac{5x^{4} - 34x^{3} + 70x^{2} - 33x - 19}{(x - 3)^{2}}$

\begin{solution}
$\frac{5x^{4} - 34x^{3} + 70x^{2} - 33x - 19}{(x - 3)^{2}} = 5x^{2} - 4x + 1 + \frac{9}{x - 3} - \frac{1}{(x - 3)^{2}}$
\end{solution}
\end{problem}

\begin{problem}
$\frac{x^{6} + 5x^{5} + 16x^{4} + 80x^{3} - 2x^{2} + 6x - 43}{x^{3} + 5x^{2} + 16x + 80}$

\begin{solution}
$\frac{x^{6} + 5x^{5} + 16x^{4} + 80x^{3} - 2x^{2} + 6x - 43}{x^{3} + 5x^{2} + 16x + 80} = x^{3} + \frac{x + 1}{x^{2} + 16} - \frac{3}{x + 5}$
\end{solution}
\end{problem}

\begin{problem}
$\frac{-7z^{2} - 76z - 208}{z^{3} + 18z^{2} + 108z + 216}$

\begin{solution}
$\frac{-7z^{2} - 76z - 208}{z^{3} + 18z^{2} + 108z + 216} = -\frac{7}{z + 6} + \frac{8}{(z + 6)^{2}} - \frac{4}{(z + 6)^{3}}$
\end{solution}
\end{problem}

\begin{problem}
$\frac{-10z^{4} + z^{3} - 19z^{2} + z - 10}{z^{5} + 2z^{3} + z}$

\begin{solution}
$\frac{-10z^{4} + z^{3} - 19z^{2} + z - 10}{z^{5} + 2z^{3} + z} = -\frac{10}{z} + \frac{1}{z^{2} + 1} + \frac{z}{(z^{2} + 1)^{2}}$
\end{solution}
\end{problem}

\begin{problem}
$\frac{4s^{3} - 9s^{2} + 12s + 12}{s^{4} - 4s^{3} + 8s^{2} - 16s + 16}= \frac{\answer{1}}{s - 2} + \frac{\answer{4}}{(s - 2)^{2}} + \frac{\answer{3s + 1}}{s^{2} + 4}$
\end{problem}

\begin{problem}\label{findparfraclast}
$\frac{2s^{2} + 3s + 14}{(s^{2} + 2s + 9)(s^{2} + s + 5)}$

\begin{solution}
$\frac{2s^{2} + 3s + 14}{(s^{2} + 2s + 9)(s^{2} + s + 5)} = \frac{1}{s^{2} + 2s + 9} + \frac{1}{s^{2} + s + 5}$
\end{solution}
\end{problem}

\end{question}

\begin{problem}
Find a partial fraction decomposition of $R(z) = \frac{4}{z^4-1}$ over the \textit{complex} numbers.

\begin{solution}
$R(z) = \frac{4}{z^4-1} = \frac{1}{z-1} - \frac{1}{z+1} + \frac{i}{z-i} -  \frac{i}{z+i} $
\end{solution}
\end{problem}

\begin{problem}
In light of Theorem \ref{complexfactorization}, we know all polynomial functions can be reduced to a product of \textit{linear} factors - if we use complex numbers.  It turns out that Theorem \ref{pfdecomp} holds true with complex numbers as well (though when using complex numbers, there are no irreducible quadratics.)  Discuss with your classmates how this ultimately means every rational function is a sum of shifted Laurent Monomials.\footnote{See Section \ref{IntroRational}.}
\end{problem}  

\begin{problem}
One of the most common algebraic error the authors encounter when teaching  Calculus is along the lines of:

\[ \frac{8}{x^2 - 9} \neq \frac{8}{x^2} - \frac{8}{9}\]

Think about  why if the above were true, this section would have no need to exist.
\end{problem}


\end{document}
